\section{Hub Siting and Coverage Analysis}

\subsection{T1/T1 Diamond Hub Locations}

The nine T1/T1 diamond hub locations are determined by the intersection of T1 Primary Artery corridors, as identified in the ROUTE Track A and B analysis \citep{ROUTE_A1,ROUTE_B4}. These are not chosen for transit suitability; they are the freight network's primary node locations, selected because they represent the highest-importance intersections of the highest-volume corridors. The transit suitability analysis presented here is performed on the pre-determined freight sites.

Table~\ref{tab:hubs} presents the nine hub locations with transit-shed population estimates. Population figures are derived from the 2022 ACS 5-year estimates for all census tracts with centroids within 30 miles of the hub location \citep{ACS_B08201_2022}. Transit-dependent household counts are ACS B08201 zero-vehicle household estimates for the same geography.

\begin{table}[h]
\centering
\caption{T1/T1 Diamond Hub Locations: Transit-Shed Analysis (30-mile radius)}
\label{tab:hubs}
\begin{tabular}{llrr}
\toprule
Hub Location & Corridors & Population (30mi) & Transit-Dependent HH \\
\midrule
Atlanta, GA & I-75/I-85 & 5,800,000 & 180,000 \\
Gary/Chicago, IL & I-80/I-90 & 8,100,000 & 320,000 \\
Boston, MA & I-95/I-90 & 4,200,000 & 115,000$^*$ \\
Seattle, WA & I-5/I-90 & 3,800,000 & 115,000 \\
Sacramento, CA & I-5/I-80 & 2,400,000 & 75,000 \\
San Antonio, TX & I-10/I-35 & 2,500,000 & 88,000 \\
Richmond, VA & I-95/I-85 & 1,300,000 & 42,000 \\
Jacksonville, FL & I-10/I-95 & 1,400,000 & 48,000 \\
Toledo, OH & I-75/I-90 & 700,000 & 28,000 \\
\midrule
\textbf{Total (non-overlapping)} & & \textbf{26,800,000} & \textbf{896,000} \\
\bottomrule
\end{tabular}
\end{table}

\noindent $^*$Boston has strong existing transit; hub value is primarily for intercity connections to underserved New England communities.

The total 30-mile population of 26.8 million overstates the non-duplicating coverage because several hub catchment areas overlap: Gary/Chicago and Toledo share population in the Lake Michigan--Lake Erie corridor, and Atlanta and Jacksonville share population in north Florida. After overlap correction using population centroid assignment to the nearest hub, the non-overlapping transit-shed population is 22.1 million total residents, of whom an estimated 12.4 million are transit-dependent (zero-vehicle households or rural households without viable alternative transportation).

\subsection{T1/T2 Regional Stop Identification}

In addition to the 9 T1/T1 diamond hubs, the I2.0 program identifies approximately 50 T1/T2 regional stops: locations where a T1 Primary Artery intersects a T2 Major Connector, creating a natural pull-off location for intercity bus service to smaller communities within 20 miles of the interchange.

T1/T2 stop selection criteria:
\begin{enumerate}
  \item Existing rest area within 0.5 miles of the T1/T2 interchange
  \item Population center of 5,000 or more residents within 20 miles of the interchange
  \item No existing Greyhound or FlixBus stop within 25 miles
  \item T2 corridor intersection with ADT $\geq$ 15,000 (indicating viable access road traffic)
\end{enumerate}

Representative T1/T2 stops include:
\begin{itemize}
  \item \textbf{Meridian, MS} (I-20/US-45): population within 20 miles = 68,000; no current bus service following 2020 Greyhound restructuring
  \item \textbf{Abilene, TX} (I-20/US-83): population within 20 miles = 127,000; nearest Greyhound terminal 90 miles
  \item \textbf{Salina, KS} (I-70/I-135): population within 20 miles = 58,000; regional hub for central Kansas with no intercity bus connection
  \item \textbf{Twin Falls, ID} (I-84/US-93): population within 20 miles = 81,000; no rail, nearest bus terminal 195 miles (Boise)
  \item \textbf{Amarillo, TX} (I-40/US-87 junction): population within 20 miles = 213,000; existing Greyhound terminal but no northbound service toward Lubbock/Abilene
\end{itemize}

\subsection{Transit-Shed Radius Analysis}

The choice of 30-mile catchment radius as the primary analysis unit reflects the tradeoff between transit-shed coverage and realistic access behavior. A 20-mile radius is appropriate for urban hubs where bus or ride-share access is available; a 30-mile radius reflects the rural context in which most T1/T2 stops operate, where drive-to-hub access is the primary mode \citep{Pisarski2006}. A 50-mile radius produces substantially larger population estimates but likely overstates realistic travel-to-hub behavior for zero-vehicle households, who by definition cannot drive themselves to the hub.

\begin{table}[h]
\centering
\caption{Transit-Shed Population by Radius: 9 T1/T1 Hubs (non-overlapping)}
\label{tab:radii}
\begin{tabular}{lrr}
\toprule
Radius & Total Population & Transit-Dependent HH \\
\midrule
20 miles & 18,400,000 & 7,100,000 \\
30 miles & 22,100,000 & 8,900,000 \\
50 miles & 29,600,000 & 11,200,000 \\
\midrule
T1/T2 stops (30 mi) & 4,800,000 & 3,500,000 \\
\midrule
\textbf{Combined (30 mi, non-overlap)} & \textbf{26,900,000} & \textbf{12,400,000} \\
\bottomrule
\end{tabular}
\end{table}

The 12.4 million figure used in the abstract and throughout this paper is the non-overlapping combined transit-dependent population within 30 miles of either a T1/T1 diamond hub or a T1/T2 regional stop. This is the primary estimate. For policy analysis, the more conservative 20-mile radius estimate (7.1 million from T1/T1 hubs alone) provides a floor; the 50-mile radius estimate (11.2 million from T1/T1 hubs plus 3.5 million from T1/T2 stops) provides a ceiling.

\paragraph{Proximity vs.\ Operational Access.}
The population figures in Table~\ref{tab:hubs} measure \textit{geographic proximity} ---
the number of residents living within 30 miles of a hub location --- not operational transit
access. For transit-dependent travelers (zero-vehicle households), geographic proximity to a
hub located on a controlled-access freight corridor does not constitute access without a local
feeder connection (van service, bus route, or dial-a-ride) linking the community to the hub
location. An estimated 12.4 million transit-dependent Americans live within 30 miles of a
T1/T1 hub location \textit{and would benefit from a transit connection if one were provided.}
This paper analyzes the hub increment investment (passenger platform, parking, amenities) at
the highway location; it does not include the cost of local feeder services, estimated at
\$500--800 million nationally for a baseline van/shuttle network connecting the nearest
communities to each hub (see Section~\ref{sec:marginal-cost}). The 14$\times$ efficiency
finding --- \$161 per traveler-served versus \$2{,}200 for standalone terminals --- is based
on the hub increment cost only and does not include feeder costs.

\subsection{Existing Service Gap Confirmation}

To confirm that the estimated 12.4 million transit-dependent residents lack existing intercity bus service, we cross-referenced the T1 hub 30-mile catchment areas with the BTS Intermodal Passenger Connectivity Database \citep{BTS_intermodal2023} and the current Greyhound/FlixBus/Megabus route maps as of January 2026. The cross-reference shows:
\begin{itemize}
  \item Of the 50 T1/T2 stops, 43 have no intercity bus stop within 25 miles
  \item Of the 9 T1/T1 hubs, 6 have existing bus service at the hub city but none connecting to the T1/T2 stop network for rural intermediate service
  \item The Gary/Chicago hub has the strongest existing transit, but the south side communities within 30 miles have below-average access to Chicago's existing intercity terminals
\end{itemize}
